\chapter{Optimisation of the CWP algorithm based on the morris sensitivity analysis}

\section{Deploying insight from the morris screening}

% The results of the morris sensitivity analysis have shown that the buffer pixel seems to be largely non influential on the area of cwp. Therefore the strategy of factor fixing can be performed, and the buffer pixel parameter can be internalized and be set to a constant value which in this case was choosen to be 3 which lies within the range of the morris screening grid. 

% Further it is evidential, that the step length seems to have a reasonable influence on the computed patch area both for the tributary caused and non-tributary caused cwp. It is therefore of interest to remove the dependency of the algorithm on this parameter aswell.

% Laslty the temperature delta parameter seems to have the strongest effect of all tested parameters. This parameter however can not be simply reduced away, as it has to be set by the user. This because in litterature different temperature deltas for the definition of a cwp exist and therefore a user needs to be able to pick the one that is corresponding with his or heres definition.


% So summing up the new algorithm can remove the buffer pixel parameter by factor fixing and has to provide a new approach of detecting ecologically reasonable cwp without relying on a step length parameter.


\section{Operationalisation of Surrounding Stream temperature}

% The later was achieved by rethinking the operationalsiation of the idea of surrounding stream temperature as given by the ecological definition. Instead of relying on the mean temperature of fixed slabs, the new algorithm deploys a inverse distance weighting scheme to compute a reference temperature for each cell by deploying a inverse distance weighting scheme. This guarantees that the reference temperature for pixel flagging is always a mutliperspective decission taking surrounding temerature values of different locations into account to provide a surrounding temperature. This strategy as the added benefit that the deployed reference temperature is a spatially continuous phenomenon unlike the older purely slab based approach. Therefore sharp cutoffs of cwps at slab boundaries can no longer occur unlike in the old algorithm 

% => here an image of this sharp cut of


\section{Implementation Detail}

% generally the implementation to a large degree follows the implementation of the time optimized cwp algorithm only differing from said function in the computation of the reference temperature step.

% the following pseudo code provides a detailed overiev of this new reference temperature computation


% 1. From a uniform distribution with bounds 500 3000 sample 10 values as step length
% 2. For each of the 10 step length values:
	% devide the river reach into n centerline segments based on the step length parameter.
	% construct a set of 3 pixel wide narow reference strip polygon for each slab.
	% round the TIR raster using a fixed rounding factor.
	% extract the median temperature of each segment from the rounded tir raster using the refference strip polygon.
	% Get the "visual center of each polygon" as a point and assign the computed reference temperature to this point
% 3. Initiate a new raster which will be populated by reference temperature values
% 4.  For each cell in this new raster:
	% if the cell at the same location in the TIR raster has an NaN value:
		% set the pixel to nan and continue to the next pixel
	% else
		% compute a weighted temperature value for this pixel using inverse distance weighting of the reference points.
% 5. Subtract the rounded TIR raster from the reference raster; flag pixels where the difference exceeds dlta T producing a binary CWP raster.
% 7. polygonize the binary CWP raster into patch polygons
% 8. Filter patches by minimum area, compute per-patch temperature statistics and apply a patch-level threshold check against the reference temperature.


% Step 4. is thereby written in C++ and imported into R using the package Rcpp.


